\paragraph{The cusp test on the principal marked charts.}
Retain the period chart and identity Fourier mark of
Propositions~\ref{bfam:unitread} and \ref{bfam:oncharts}.
The following argument uses the mark at one prime over \(M\).
In particular it applies before summing over the definite column genus.

\begin{lemma}\label{crev:mode}
For every rational isotropic row line and every orthonormal frame on
the chart, the real diagonal transvection restriction of each marked
theta jet has mean zero along any period with trivial multiplier.
Its full real cusp-torus average is zero as well.
\end{lemma}
\begin{proof}
Write \(v\ne0\) for the isotropic row and \(I=(Ev)^n\) in the
rational trace symplectic space of the oscillator. Choose a rational
Lagrangian in one row containing \(Ev\), and take its \(n\)-fold sum
\(D\). Complete it to a symplectic polarization \(D\oplus D^*\),
using \(D^*\) for the position variable. Translations by \(a\in
D(\mathbb A_f)\) are multiplication by the finite additive character
of \(z(a)\) at a position \(z\). In particular they act as the identity
at all \(z\) annihilating \(I\).

Let \(\varphi\) be the marked finite vector in these coordinates.
The finite translation calculation preceding Lemma~\ref{bfam:genus}
gives an element of \(I(\mathbb A_f)\) acting on the single marked
vacuum with eigenvalue different from one. Indeed, at a prime of
\(\mathfrak m\) its first \(r\) residues are a basis and a primitive
vector on the conjugate line, divided by the local half-level,
detects one of them. The unshifted character of the isotropic line
is one. It follows pointwise that
\[
                 \varphi(z)=0\qquad (z|_{I(\mathbb A_f)}=0).
\]
This also holds for a sum of the indicated marks, term by term. The
calculation takes place after changing polarization, which intertwines
the Weyl translations by the Fourier and substitution laws in the
proof of Proposition~\ref{bfam:unitread}.

In this polarization the real diagonal flow is a shear. Its
infinitesimal action \(N\) has image \(I_{\mathbb R}\) and kernel
\(I_{\mathbb R}^{\perp}\), and it fixes \(D_{\mathbb R}\) pointwise.
The shear contributes the character with quadratic exponent,
up to the fixed Fourier sign, \(Q(z)=\tfrac12\langle \widetilde
z,N\widetilde z\rangle\), where \(\widetilde z\) is a lift of the
position and the pairing is symplectic. At each place the trace
polarization applied to the transvection
\(\kappa_0 h(y,v)v\) gives a nonzero real multiple of
\(\sum_{j=1}^n|h_\tau(y_j,v)|^2\). The multiples have the same sign:
they are obtained from the product of \(\tau\xi_0\) and
\(\tau\kappa_0=\pm2\tau\xi_0\) with one fixed choice of sign, and
the scaling \(k_0\) of the real parameter is positive. Consequently
\(Q(z)=0\) exactly on the annihilator of \(I_{\mathbb R}\).

More explicitly, on the real flow the first period matrix is fixed
because the shear fixes the chosen first Lagrangian. The second
matrix changes by the first times the real symmetric shear matrix.
Thus \(\Im\Omega\), the holomorphic argument \(hz_j\), the Gaussian
prefactor and \(\det h\) in Proposition~\ref{bfam:unitread} stay
fixed. The change in each summand is exactly its quadratic phase,
also after applying the polynomial in the \(z_j\)-derivatives.
The real restriction therefore has the expansion
\[
            \sum_{\ell\in D^*(\mathbb Q)}
                 \varphi(\ell)\,A_x(\ell)\,e^{2\pi i s Q(\ell)}
\]
with the Fourier sign absorbed in \(Q\). At the fixed frame the
archimedean coefficient is a polynomial times a decaying Gaussian,
and the compact finite support restricts rational positions to a
fixed denominator lattice. The series converges absolutely and
uniformly for real \(s\). If a rational position annihilates \(I\)
at infinity it annihilates \(I\) at all finite places. Hence every
zero-frequency summand vanishes. Averaging over \([0,T]\) and
taking \(T\to\infty\) gives zero by dominated convergence, even
when individual frequencies are not characters of the chosen
period circle. For a periodic sum this is its period mean.
Real translates of the starting frame give the same conclusion.
Integrating along the diagonal circles in the full cusp torus now
gives the torus average by translation invariance of Haar probability.
\end{proof}

\begin{proposition}\label{crev:half}
On each of these principal period charts the raw theta readings at
\(n=8(fr+1)\) are bounded in Petersson norm. For every nonzero
reading each orthonormal frame attaining at least half the component
supremum has row minimum squared at least \(\varepsilon\) of
Proposition~\ref{bergvol:global}. The same minimum conclusion applies
to every nonzero bounded cusp form with trivial transvection multiplier on
the chart. In particular the cusp and half-maximum tests on the raw
rules in Theorem~\ref{bfam:oncharts} follow from its stated marked
period uniformization.
\end{proposition}
\begin{proof}
Write \(a(x)=\min_{0\ne v\in L'}\sum_\tau{\rm maj}_{x,\tau}(v)>0\).
Lemma~\ref{bfam:minimum} and Proposition~\ref{intjet:arch}, at
orthonormalized frames, first give a bound of the form
\[
             \|P(x)\|_{\rm Pet}\le C(1+a(x)^{-1})^B .
\]
Indeed \(\lambda^2\ge 2a(x)/R_E^2\), the theta generator has
mean square one, and the directional norm and its weight correction
are fixed on orthonormal frames. Here \(C,B\) may depend on the
fixed section and stage. Multiplication by its nonzero fixed
reflex scalar does not affect this assertion.

Put \(H=2M\sqrt d\). For \(a=a(x)<\min(1,H^{-1})\) choose a
shortest \(v\). Its integral self-pairing vanishes by the norm
argument following Proposition~\ref{bergvol:global}. The
transvection there acts trivially modulo \(M\) on \(\mathfrak a L'\),
so lies in the actual marked stabilizer by Lemma~\ref{munf:marks}.
Its multiplier on the raw sections is trivial by the eighth-power
and identity-mark calculation in Theorem~\ref{bfam:oncharts},
which uses only the transformation formula. Take \(H\) as period
in the scaled coordinate, and write \(G(\zeta)\) for the raw
coefficient along this flow. It is holomorphic on
\(\Im\zeta>-1/a\) by the leading-minor calculation above.

We give the majorant comparison needed to apply the polynomial bound
before knowing a global supremum. At a single place set
\(l=\|\tau v\|_x^2\le a\). Rotate the positive and negative blocks
so the two components of \(\tau v\) lie on their first coordinate
axes, each of norm \(\sqrt{l/2}\). These rotations preserve the
central majorant. In this plane the new positive line has positive
and negative coordinates \(A=1-i\zeta l/2,\ B_1=-i\zeta l/2\),
with \(|A|^2-|B_1|^2=1+l\Im\zeta\). All orthogonal directions
retain their majorant. The two comparison eigenvalues in the plane
are \(K\) and \(K^{-1}\), where
\[
 K=\frac{(|A|+|B_1|)^2}{1+l\Im\zeta}
    \le\frac{(1+l|\zeta|)^2}{1+l\Im\zeta}.
\]
Indeed a unitary change of phase makes the normalized line coordinates
nonnegative and its orthogonal negative mate interchanges them.
At \(\zeta=s-i/(2a)\), \(0\le s\le H\), this gives \(K\le 13\):
the denominator is at least \(1/2\) and
\(l|\zeta|\le Ha+1/2\le3/2\).
Thus the new row minimum is at least \(a/13\).
The leading-minor variation estimate preceding
Lemma~\ref{bfam:harm} bounds the conversion from Petersson norm
to raw coefficient on this line by \(2^{w/2}\).
It follows that
\[
      \sup_{0\le s\le H}|G(s-i/(2a))|
                    \le 2^{w/2}C(1+13/a)^B .
\]
For \(y=\Im\zeta\ge0\), \(0\le s\le H\), the same comparison
gives \(K\le(1+a(H+y))^2\). The leading-minor formula bounds
the raw metric conversion by \((1+ay)^{w/2}\).
Together these give polynomial growth of \(G(s+iy)\) in \(y\)
for the fixed starting frame. In \(q=\exp(2\pi i\zeta/H)\),
Cauchy's coefficient estimates therefore remove all negative
Laurent powers. Lemma~\ref{crev:mode} removes the constant power.
The maximum principle for \(G/q\) on the disk of radius
\(\exp(\pi/(aH))\) gives
\[
      \|P(x)\|_{\rm Pet}=|G(0)|
       \le 2^{w/2}C(1+13/a)^B e^{-\pi/(aH)} .
\]
This is bounded as \(a\downarrow0\). On the complementary range
the original polynomial estimate is bounded. Hence every nonzero
section has a finite positive supremum and one can choose an
orthonormal frame above half of it. With boundedness established,
the half-maximum argument following Proposition~\ref{bergvol:global}
applies with its original \(\varepsilon\); the qualitative constants
\(B,C\) disappear. For a general bounded cusp form the real torus
and positive-frequency argument there supplies its diagonal zero
mode. Thus the same quantitative conclusion holds for that space
at every weight whose transvection character is trivial.
\end{proof}
The proof uses the short vector of the component under consideration.
It gives the absolute-constant bound of Corollary~\ref{bfam:metric}
for all selected parameters, all permitted column lattices and
auxiliary integers at every good embedding; it requires neither
a common fundamental domain as the rank changes nor a uniform
constant in the preliminary polynomial growth estimate.
\src{src/Development/AbelianVarietyTheory/Proofs/CuspFormEvaluations/ExistsExpansionOfEqSumThFr.lean}{26}{The source's different zero-mode test using positivity of the last row weight};
\src{src/Development/AbelianVarietyTheory/Proofs/DivisorProofs/FirstMinAtHalfMax.lean}{111}{Transfer to the minimum at the selected level}.

\begin{corollary}\label{crev:twist}
On the good marked charts the algebraic rule space defined by
testing all good complex readings requires no additional cusp
comparison for the raw spanning argument. A regular unit section on
a component with constant positive Petersson norm, in mutually dual
multiplicative metrics for it and its inverse, preserves this cusp
condition in both directions. This applies to the key and the
horizontal determinant whenever they exist on that component with
their stated constant metrics.
\end{corollary}
\begin{proof}
In the rule definition the cusp subspace is the intersection of
the regular weight rules with the inverse images, under every good
reading, of the complex spans of the analytic cusp forms. Each such
span is the analytic cusp space itself: holomorphy, transformation
law and vanishing of the torus integrals are linear, and each
finite sum of bounded forms is bounded. Conversely inclusion in the
span holds for each member. Lemma~\ref{bfam:rule},
Lemma~\ref{crev:mode} and Proposition~\ref{crev:half} place the
raw rules in this intersection. For generation one only needs to
express the reading of every member in the span of the raw readings,
which is the conclusion of Proposition~\ref{munf:span} on those
charts. Injectivity after scalar extension in Lemma~\ref{bfam:rule}
then descends the span to the coefficient field.

For the multiplication assertion use a common finite-index
transvection period lattice killing both characters. Passing to
it leaves each normalized constant term unchanged when applied to
a function periodic for the original lattice. Along a rational
isotropic row the positive-ray metric estimate following
Proposition~\ref{bergvol:global} and boundedness imply that
every nonzero Fourier coefficient of either factor has all
coordinates nonnegative, by growing one imaginary coordinate
at a time. At height zero the coefficients are absolutely summable
by smoothness on the real torus. Thus the product constant is the
product of the two constants, by convolution: two vectors with
nonnegative coordinates can sum to zero only if both are zero.
This proves cusp vanishing for the product, and boundedness follows
by multiplicativity of the metric. Apply the same argument to the
inverse section, which is holomorphic and has the reciprocal
constant norm. In particular its Fourier support is tested afresh
using holomorphy in the same upper tube, even though its \emph{weight}
has the opposite sign. The polynomial growth bound uses only the
weight spread. This treats the Fourier condition in the rule space;
a separate use of a subcanonical sheaf on a compactification must
identify its restriction with that space.
\end{proof}
\src{src/Formalism/LocalToGlobal/ProjectiveLattices/CuspFree.lean}{14}{All good embeddings in the definition of the cusp rules};
\src{src/Formalism/LocalToGlobal/ProjectiveLattices/JFDFree.lean}{42}{The span of the bounded holomorphic cusp readings}.

\paragraph{The flag measure of the selected unitary genus.}
Let \(F_0=E^+\). Use the host fields of Proposition~\ref{host:art},
the genus of the self-dual row, and its odd principal level \(M\ge3\).
Write \(K\) for the finite adelic stabilizer and \(K(M)\) for its
principal subgroup. Let \(c_r\) and \({\rm mass}_I\) be as in
Proposition~\ref{bergvol:compactmass}.

\begin{proposition}\label{crev:mass}
For the full adelic packet of marked rows, the sum of the component
flag volumes in the normalization of Lemma~\ref{bergvol:local} is
\[
        V=[K:K(M)]\,\frac{c_r^f\,{\rm mass}_I}{(2\pi)^{fr}}
           \le V_* .
\]
Adelic probability on flag-invariant functions pushes forward to
the sum of the fundamental-set measures divided by \(V\).
Conditioning on a nonempty union of components gives the same
formula with its own total flag volume.
\end{proposition}
\begin{proof}
In the marked period description of Lemma~\ref{munf:marks}, the
finite double cosets of \(G=U(h_0)\) represent the marked rows of
this genus. Indeed finite translates of the completed reference
lattice and its basis give a row by intersection with the rational
space. Conversely choose local isometries to a specified row in
the genus. At primes of \(M\) the extension is split and the mark
is a basis on one half with its pairing-dual on the other. An
invertible matrix over the local residue quotient lifts to an
integral invertible matrix by lifting entries. Thus the local
isometries may match the marks and the stabilizer is \(K(M)\).
In terms of the two algebraic Lie ranks \(p,q\), the real flag
space at a place is \(U(p_0,r-p_0)/U(1)^r\), with \(p_0=q\).

Use the rational form \(J=\operatorname{diag}(1_{p_0},-1_{r-p_0})\) given
by Lemma~\ref{cmeig:row}, and compare with \(G_c=U(1_r)\), the
group of the positive definite standard form. At every finite
place the two self-dual lattices are isometric, including residue
characteristic two, by the local part of that lemma. Keep track of
the Haar densities in this assertion. On the Lie algebra for a
Hermitian matrix \(J\), left translation identifies an invariant
top differential by identifying its tangent at one with the fixed
\(F_0\)-space of skew-Hermitian matrices via \(X\mapsto JX\).
Choose the same rational top differential on that vector space
for both \(J\) and \(1_r\). Under an isometry \(P^*JP=1_r\),
the comparison in these coordinates is \(A\mapsto P^*AP\).
Its determinant on the skew-Hermitian space is
\(N_{E/F_0}(\det P)^r\): extend scalars to the split algebra,
where the two factors multiply an unrestricted \(r\)-by-\(r\)
matrix on its left and right. Its absolute value is one at each
finite place because \(N(\det P)=1/\det J=\pm1\).
Thus the corresponding maximal stabilizers have the same local
volume. No choice of their local bases contributes a discriminant.

The global scalar and convergence factors in the Tamagawa measure
are the same: the base field and dimension agree and the geometric
character lattice in both cases is generated by determinant with
the quadratic Galois action of \(E/F_0\). At each real place the
common gauge in the \(A\)-coordinates is the same scalar times
Lebesgue measure on the imaginary diagonal and the complex
strictly lower entries. Multiplication by \(J\) only changes signs
in these coordinates. Divide out that same scalar for both groups.
Use volume one for the maximal stabilizers at finite places and
the resulting real coordinate densities. In these choices the
total adelic covolumes of the two groups are equal: they are
obtained from Tamagawa measure by a common scalar and both
Tamagawa numbers are two for a nonzero Hermitian space over a
quadratic extension. See \href{https://doi.org/10.1007/BF02391774}
{Weil, \emph{Sur la formule de Siegel dans la th\'eorie des groupes
classiques}}, and Ichino, \emph{A regularized Siegel--Weil formula
for unitary groups}, \S4. For the definite group the real volume is
\(c_r^f\). The finite double cosets at maximal level are its
standard lattice genus: cutting the quotient over such a coset
divides this common real volume by its finite automorphism group.
Summing gives \(c_r^f\,{\rm mass}_I\).

For the indefinite group, the quotient of the real density by
angle measure \(d\theta_1\cdots d\theta_r\) is the stated flag
measure. At the central flag the tangent projection takes the
strictly lower entries of the Lie algebra with Jacobian one.
At a lower unipotent matrix \(u=k b\), with \(k\) unitary and
\(b\) upper triangular, the projected tangent in this frame is
the strictly lower part of \(b(u^{-1}du)b^{-1}\).
The unipotent differential has determinant one in its entries.
Conjugation on matrices modulo upper triangular matrices has
complex determinant \(\prod_{i>j} b_{ii}/b_{jj}\). Its squared
absolute value is \(\prod_{i=1}^{r-1}D_i^{-2}\), using
\(|b_{ii}|^2=D_i/D_{i-1}\) and \(D_0=D_r=1\).
This recovers the chart density at every point.

Over each finite double coset \(gK(M)\), the quotient of
\(G_\infty\times gK(M)\) is by
\(\Gamma_g=G(F_0)\cap gK(M)g^{-1}\). Cutting the real factor
by a fundamental set for \(\Gamma_g\) gives the common finite
volume \(1/[K:K(M)]\) times real quotient volume. There is no
stabilizer on the flag: it would be discrete in a compact torus,
hence finite, and the principal subgroup on the underlying
integral period lattice is torsion-free. For the latter, at an
odd prime dividing \(M\) the powers of a nonidentity
\(1+\ell^a B\), \(a\ge1,\ B\not\equiv0\bmod\ell\), have depth
\(a+v_\ell(t)\) at exponent \(t>0\), by the binomial theorem.
Integrating along the angle tori therefore gives \((2\pi)^{fr}\)
times the flag integral. The same disintegration applies to any
nonnegative measurable function of the flag. Applying it first
to one proves the identity and then dividing by the total mass
proves the probability formula. The congruence index is at most
\(M^{2fr^2}\) by reduction on the row, so
Proposition~\ref{bergvol:compactmass} gives the required bound.
\end{proof}
In particular both analytic tests in
Proposition~\ref{bergvol:global} are available on the selected
principal packet for the bounded analytic cusp space of trivial
transvection character. Using this for an algebraic component
uses its identified framed period charts. For the rule space the
reading is as in Corollary~\ref{crev:twist}; agreement with a
separately defined toroidal sheaf is a geometric comparison.

\begin{corollary}[Analytic estimates on a selected good cover]\label{crev:good}
Take a stage in the selected range on the host of
Proposition~\ref{host:art}, a core with the prescribed host,
polarized datum and level, and the deep cover \(Y\) with
\(Y.\mathrm{Good}(e,o)\). Use its identity auxiliary rule
at \(n=8(fr+1)\), and the parameters and semiabelian tests of
Theorem~\ref{bfam:oncharts}. At every good embedding the raw
rules are cusp rules on this \(Y\). The short integral rules in
that theorem span the untwisted space on the cover and on each
component, with logarithmic Petersson norms at most
\(C(frL)^2S+2\log D_M\) for the stated stage size bounds.
For the component cusp spaces at the frame weights of
Corollary~\ref{bergvol:size} the normalized evaluation kernel
is at most \(\beta_*\).
\end{corollary}
\begin{proof}
Here \(Y.\mathrm{Good}(e,o)\) contains the following data. Its
\code{unf} gives the complex framed period uniformization at each
good embedding, including the section transformation law and the
Haar mean-square-one normalization of the theta generator on unit
frames. In these coordinates the first \(p_0=q\) Hermitian
directions are positive. Their conjugate, reversed tangent block
has Lie rank \(q\), and the other tangent block has Lie rank \(p\)
at the labelled \(\tau\). Thus the pivot \(\lambda_{\tau,p}\)
in Theorem~\ref{bfam:oncharts} is numbered in the algebraic
weights, while the positive block used for the majorant ends at
\(p_0\) in the Hermitian flag.

The \code{modelW} and \code{xiAbs} clauses give the signed host
Gram matrix and the absolute polarization coefficient at every
place. The common lattice embedding and alternating pairing in
\code{adPin} identify the original period lattice, its finite
adelic translates, the half-level mark and the eigenframe
coordinates on these same components. In particular stripping the
fixed host ideal from a translate of the reference period
\(\mathfrak a L'_*\) gives the translated Hermitian self-dual row:
finite unitary matrices commute with \(\mathfrak a\), and
localization detects self-duality. This is \(L'\) in the return
count and in Lemma~\ref{bfam:minimum}. The \code{lev} clause gives
auxiliary integer one on the deep variety; the spanning construction
there uses its identity auxiliary map of Corollary~\ref{bfam:mark}.
It does not replace the original map from the deep variety to the
base by an identity.

The \code{dictU} and \code{levelK} clauses identify the rational
and right-principal-level equivalences with the chart action and
exhaust the real flag orbits. Indeed the adelic frame map is
equivariant modulo arithmetic translates under the rational and
level actions and intertwines the real action; it hits each flag
modulo its upper triangular frame change. Equality of its frames
modulo the arithmetic action forces rational and finite-level
equivalence. The unit frames over each flag form a single orbit
under its diagonal angle group. Together with the marked-period
classification in Lemma~\ref{munf:marks}, this is the packet and
the probability disintegration of Proposition~\ref{crev:mass}.
The \code{unit} clause gives leading Gram determinants of absolute
value one at the adelic frames, as used in their metric reading.

Finally \code{chiTriv} gives trivial arithmetic multiplier at the
dominant natural weights on each good chart. It gives the same
for every integral flag weight: for fixed arithmetic element its
multiplier is multiplicative in addition of weights and never
zero. Take natural weights with \(n=0\) that are one on an
initial segment at one labelled embedding and zero elsewhere.
Consecutive differences give all individual weight coordinates,
with reversed order and a minus sign on the dual block. Their
integral span is the whole flag weight lattice. In particular
there is no character loss in the transvection period for the
frame weights. Apply Propositions~\ref{crev:half} and
\ref{crev:mass}, then Proposition~\ref{bergvol:global}.
The actual all-weight space \code{cuspW} is the intersection of
the regular rules with the good-reading tests of
Corollary~\ref{crev:twist}; at a natural weight it is the ordinary
rule space \code{cusp}, and \code{cuspOn} adds only the component
support condition. Thus Theorem~\ref{bfam:oncharts} applies with
its integral tests and no additional analytic cusp assumption.
\end{proof}
\src{src/Development/AbelianVarietyTheory/Definitions/Divisors/NormNegCmConstantEqAtAllEmbeddings.lean}{187}{The clauses of the selected deep cover};
\src{src/Formalism/LocalToGlobal/ProjectiveLattices/FamilyUniformisationBasePointOrthogonalBis.lean}{97}{The marked finite translates in the period identification};
\src{src/Development/AbelianVarietyTheory/Definitions/Foundations/IsProbabilityMeasureMuPart.lean}{425}{The actual all-weight rule space};
\src{src/Development/AbelianVarietyTheory/Counterexamples/LevelGroupExamples/AdelEqZeroOfOffDom.lean}{77}{Generation of the arithmetic characters by natural weights}.

Constructing such covers at all scales is a separate existence
statement. For example, the sized-cover construction starts with
the pinned core, a stage in range, the host splitting bound at least
five and \(15\) coprime to the base level. It constructs the half-level,
uses PEL moduli representability, forms the coefficient compositum
with a ray class field, and builds the marked uniformization; the
field-size bound uses the ray-class-field estimate. If the comparison
to a prescribed base frame requires \(Y.\mathrm{Match}\), that
condition concerns the same base embedding and base frames and is
not part of \(\mathrm{Good}\). The corollary asserts the estimates
on the selected \(Y\); it does not construct the core, the compatible
family or its toroidal compactifications.
\src{src/Development/AbelianVarietyTheory/Proofs/DivisorProofs/ExistsGoodCoverOverSizedField.lean}{111}{Inputs and construction of the sized good level};
\src{src/Development/AbelianVarietyTheory/Definitions/Divisors/NormNegCmConstantEqAtAllEmbeddings.lean}{207}{The common marked frames included in goodness}.
